% GENERATED FILE. Edit canonical lessons, then run npm run generate:books.
% canonical-source-hash: fnv1a64:33b2125b6e6c86d1
% canonical-lessons: ML-C408-pariganikkuka, ML-C408-vilayiruttuka, ML-C408-nilapatu, ML-C408-manobhavam, ML-C408-nyayikarikkuka

\chapter{Weighing It Up}
\label{ch:ml-weighing-it-up}
\hlchaptermodality{408}

\begin{chapteropening}
\textbf{By the end of this chapter:} \emph{I can say to consider, to evaluate, a stance, an attitude and to justify.}

Complete the last lesson of chapter 408: I can say to consider, to evaluate, a stance, an attitude and to justify.
\end{chapteropening}

\section[\texorpdfstring{\ml{പരിഗണിക്കുക} (parigaṇikkuka)}{പരിഗണിക്കുക (parigaṇikkuka)}]{\ml{പരിഗണിക്കുക} (parigaṇikkuka) — to take into account; to consider}

\label{lesson:ML-C408-pariganikkuka}



\begin{quote}
Before the new one: say the Malayalam for reality, then the Malayalam for to imply.
\end{quote}

\subsection*{You'll want to know: \ml{പരിഗണിക്കുക}}
\textbf{\ml{പരിഗണിക്കുക}} — \emph{parigaṇikkuka} — \textquotedblleft{}to take into account; to consider\textquotedblright{}.

\begin{quote}
\ml{കമ്പനി} \ml{എന്റെ} \ml{അപേക്ഷ} \ml{പരിഗണിച്ചു}.
\end{quote}

\begin{quote}
\emph{kampani enṟe apēkṣa parigaṇiccu.}
\end{quote}

\begin{quote}
The company considered my application.
\end{quote}

Past: \ml{പരിഗണിച്ചു}, \emph{parigaṇiccu}.

\ml{അപേക്ഷ} \ml{പരിഗണിക്കുക}: what is considered takes -\ml{നെ} or stays unmarked.

\begin{cousinweb}
\textbf{\ml{പരിഗണിക്കുക}} (\emph{parigaṇikkuka}) is built from Sanskrit \emph{pari-}, \textquotedblleft{}around, all over\textquotedblright{}, and the root \emph{gaṇ}, \textquotedblleft{}to count\textquotedblright{} (\emph{gaṇita}, \textquotedblleft{}counted\textquotedblright{}, CDIAL 3996), with the Malayalam verb ending -\ml{ിക്കുക}: to consider is to count a thing over.
\end{cousinweb}

\begin{grammarlens}[title={What you've built}]
One more word for giving reasons and opinions.
\end{grammarlens}

\subsection*{Your turn}
\begin{itemize}
\raggedright
  \item \emph{Say it:} \emph{parigaṇikkuka}
  \item \emph{Say it:} \emph{parigaṇikkuka}, once more
  \item \emph{Say it:} \emph{kampani enṟe apēkṣa parigaṇiccu.} — The company considered my application.
  \item \emph{Recall:} say the Malayalam for reality, then the Malayalam for to imply, then say \emph{parigaṇikkuka} again
\end{itemize}

\begin{tcolorbox}[breakable,colback=teal!4,colframe=teal!35!black,title={Before you move on}]
What does \ml{പരിഗണിക്കുക} mean? (\textquotedblleft{}To take into account; to consider\textquotedblright{}.) Say it once more.
\end{tcolorbox}
\section[\texorpdfstring{\ml{വിലയിരുത്തുക} (vilayiruttuka)}{വിലയിരുത്തുക (vilayiruttuka)}]{\ml{വിലയിരുത്തുക} (vilayiruttuka) — to evaluate; to assess}

\label{lesson:ML-C408-vilayiruttuka}



\begin{quote}
Before the new one: say the Malayalam for to imply, then the Malayalam for to take into account.
\end{quote}

\subsection*{You'll want to know: \ml{വിലയിരുത്തുക}}
\textbf{\ml{വിലയിരുത്തുക}} — \emph{vilayiruttuka} — \textquotedblleft{}to evaluate; to assess\textquotedblright{}.

\begin{quote}
\ml{അധ്യാപകൻ} \ml{കുട്ടികളുടെ} \ml{ഉത്തരങ്ങൾ} \ml{വിലയിരുത്തി}.
\end{quote}

\begin{quote}
\emph{adhyāpakan kuṭṭikaḷuṭe uttaraṅṅaḷ vilayirutti.}
\end{quote}

\begin{quote}
The teacher evaluated the children's answers.
\end{quote}

Past: \ml{വിലയിരുത്തി}, \emph{vilayirutti}.

The noun is \ml{വിലയിരുത്തൽ}, an assessment.

\begin{cousinweb}
\textbf{\ml{വിലയിരുത്തുക}} (\emph{vilayiruttuka}) is a Malayalam compound of native Dravidian parts: \ml{വില}, \textquotedblleft{}price, value\textquotedblright{} (DEDR 5421), and \ml{ഇരുത്തുക}, \textquotedblleft{}to seat, to settle\textquotedblright{} (DEDR 480). To evaluate is to settle what a thing is worth.
\end{cousinweb}

\begin{grammarlens}[title={What you've built}]
One more word for giving reasons and opinions.
\end{grammarlens}

\subsection*{Your turn}
\begin{itemize}
\raggedright
  \item \emph{Say it:} \emph{vilayiruttuka}
  \item \emph{Say it:} \emph{vilayiruttuka}, once more
  \item \emph{Say it:} \emph{adhyāpakan kuṭṭikaḷuṭe uttaraṅṅaḷ vilayirutti.} — The teacher evaluated the children's answers.
  \item \emph{Recall:} say the Malayalam for to imply, then the Malayalam for to take into account, then say \emph{vilayiruttuka} again
\end{itemize}

\begin{tcolorbox}[breakable,colback=teal!4,colframe=teal!35!black,title={Before you move on}]
What does \ml{വിലയിരുത്തുക} mean? (\textquotedblleft{}To evaluate; to assess\textquotedblright{}.) Say it once more.
\end{tcolorbox}
\section[\texorpdfstring{\ml{നിലപാട്} (nilapāṭŭ)}{നിലപാട് (nilapāṭŭ)}]{\ml{നിലപാട്} (nilapāṭŭ) — a stance; a position}

\label{lesson:ML-C408-nilapatu}



\begin{quote}
Before the new one: say the Malayalam for to take into account, then the Malayalam for to evaluate.
\end{quote}

\subsection*{You'll want to know: \ml{നിലപാട്}}
\textbf{\ml{നിലപാട്}} — \emph{nilapāṭŭ} — \textquotedblleft{}a stance; a position\textquotedblright{}.

\begin{quote}
\ml{ആ} \ml{വിഷയത്തിൽ} \ml{അവളുടെ} \ml{നിലപാട്} \ml{മാറിയില്ല}.
\end{quote}

\begin{quote}
\emph{ā viṣayattil avaḷuṭe nilapāṭŭ māṟiyilla.}
\end{quote}

\begin{quote}
Her stance on that subject did not change.
\end{quote}

To take a stance is \ml{നിലപാട്} \ml{എടുക്കുക}.

\begin{cousinweb}
\textbf{\ml{നിലപാട്}} (\emph{nilapāṭŭ}) is a Malayalam compound of native Dravidian parts: \ml{നില}, \textquotedblleft{}a standing, a place, a position\textquotedblright{} (DEDR 3675), and \ml{പാട്}, which DEDR lists with the senses \textquotedblleft{}situation, state\textquotedblright{} under the verb \ml{പെടുക} (DEDR 3853). A stance is where one stands.
\end{cousinweb}

\begin{grammarlens}[title={What you've built}]
One more word for giving reasons and opinions.
\end{grammarlens}

\subsection*{Your turn}
\begin{itemize}
\raggedright
  \item \emph{Say it:} \emph{nilapāṭŭ}
  \item \emph{Say it:} \emph{nilapāṭŭ}, once more
  \item \emph{Say it:} \emph{ā viṣayattil avaḷuṭe nilapāṭŭ māṟiyilla.} — Her stance on that subject did not change.
  \item \emph{Recall:} say the Malayalam for to take into account, then the Malayalam for to evaluate, then say \emph{nilapāṭŭ} again
\end{itemize}

\begin{tcolorbox}[breakable,colback=teal!4,colframe=teal!35!black,title={Before you move on}]
What does \ml{നിലപാട്} mean? (\textquotedblleft{}A stance; a position\textquotedblright{}.) Say it once more.
\end{tcolorbox}
\section[\texorpdfstring{\ml{മനോഭാവം} (manōbhāvaṁ)}{മനോഭാവം (manōbhāvaṁ)}]{\ml{മനോഭാവം} (manōbhāvaṁ) — an attitude}

\label{lesson:ML-C408-manobhavam}



\begin{quote}
Before the new one: say the Malayalam for to evaluate, then the Malayalam for a stance.
\end{quote}

\subsection*{You'll want to know: \ml{മനോഭാവം}}
\textbf{\ml{മനോഭാവം}} — \emph{manōbhāvaṁ} — \textquotedblleft{}an attitude\textquotedblright{}.

\begin{quote}
\ml{അവന്റെ} \ml{മനോഭാവം} \ml{എനിക്ക്} \ml{ഇഷ്ടമല്ല}.
\end{quote}

\begin{quote}
\emph{avanṟe manōbhāvaṁ enikkŭ iṣṭamalla.}
\end{quote}

\begin{quote}
I don't like his attitude.
\end{quote}

\ml{മനോ}- is the form \ml{മനസ്സ്}, the mind, takes at the front of a compound.

\begin{cousinweb}
\textbf{\ml{മനോഭാവം}} (\emph{manōbhāvaṁ}) is borrowed from a Sanskrit compound of \emph{manas}, \textquotedblleft{}mind\textquotedblright{} (CDIAL 9822), and \emph{bhāva}, \textquotedblleft{}state, disposition\textquotedblright{} (CDIAL 9475): a state of mind.
\end{cousinweb}

\begin{grammarlens}[title={What you've built}]
One more word for giving reasons and opinions.
\end{grammarlens}

\subsection*{Your turn}
\begin{itemize}
\raggedright
  \item \emph{Say it:} \emph{manōbhāvaṁ}
  \item \emph{Say it:} \emph{manōbhāvaṁ}, once more
  \item \emph{Say it:} \emph{avanṟe manōbhāvaṁ enikkŭ iṣṭamalla.} — I don't like his attitude.
  \item \emph{Recall:} say the Malayalam for to evaluate, then the Malayalam for a stance, then say \emph{manōbhāvaṁ} again
\end{itemize}

\begin{tcolorbox}[breakable,colback=teal!4,colframe=teal!35!black,title={Before you move on}]
What does \ml{മനോഭാവം} mean? (\textquotedblleft{}An attitude\textquotedblright{}.) Say it once more.
\end{tcolorbox}
\section[\texorpdfstring{\ml{ന്യായീകരിക്കുക} (nyāyīkarikkuka)}{ന്യായീകരിക്കുക (nyāyīkarikkuka)}]{\ml{ന്യായീകരിക്കുക} (nyāyīkarikkuka) — to justify}

\label{lesson:ML-C408-nyayikarikkuka}



\begin{quote}
Before the new one: say the Malayalam for a stance, then the Malayalam for an attitude.
\end{quote}

\subsection*{You'll want to know: \ml{ന്യായീകരിക്കുക}}
\textbf{\ml{ന്യായീകരിക്കുക}} — \emph{nyāyīkarikkuka} — \textquotedblleft{}to justify\textquotedblright{}.

\begin{quote}
\ml{അവൻ} \ml{തന്റെ} \ml{തീരുമാനം} \ml{ന്യായീകരിച്ചു}.
\end{quote}

\begin{quote}
\emph{avan tanṟe tīrumānaṁ nyāyīkariccu.}
\end{quote}

\begin{quote}
He justified his decision.
\end{quote}

Past: \ml{ന്യായീകരിച്ചു}, \emph{nyāyīkariccu}.

Built on \ml{ന്യായം}, fairness, a good reason, with -\ml{ീകരിക്കുക}, to make into: to make a thing look fair.

\begin{cousinweb}
\textbf{\ml{ന്യായീകരിക്കുക}} (\emph{nyāyīkarikkuka}) is built from Sanskrit \emph{nyāya}, \textquotedblleft{}rule, right method\textquotedblright{} (CDIAL 7614), with -\ml{ീകരിക്കുക}, the Sanskrit \emph{-ī-kṛ}, \textquotedblleft{}to make into\textquotedblright{}, that Malayalam uses for new verbs. It is the verb of \ml{ന്യായം}.
\end{cousinweb}

\begin{grammarlens}[title={What you've built}]
One more word for giving reasons and opinions.
\end{grammarlens}

\subsection*{Your turn}
\begin{itemize}
\raggedright
  \item \emph{Say it:} \emph{nyāyīkarikkuka}
  \item \emph{Say it:} \emph{nyāyīkarikkuka}, once more
  \item \emph{Say it:} \emph{avan tanṟe tīrumānaṁ nyāyīkariccu.} — He justified his decision.
  \item \emph{Recall:} say the Malayalam for a stance, then the Malayalam for an attitude, then say \emph{nyāyīkarikkuka} again
\end{itemize}

\begin{tcolorbox}[breakable,colback=teal!4,colframe=teal!35!black,title={Before you move on}]
What does \ml{ന്യായീകരിക്കുക} mean? (\textquotedblleft{}To justify\textquotedblright{}.) Say it once more.
\end{tcolorbox}
